\section{Every spectral derivative of the conjugate-factor defect}
\label{sd:section}

The quadratic-spectrum report \cite{QSpec} proves a first-derivative
factorization law and asks, in its second further-research question, for
the higher-derivative correction.  The answer below consists of an
ordinary convergent logarithmic sum and finitely many explicitly stated
Hurwitz or Stieltjes terms.  It works for every polynomial weight and
every positive mass.  The finite term arising from one extra power of
the spectral variable is essential.

\subsection{Conventions and the classical first derivative}

Let $a,b>0$, put $x_n=n+a$ and $Q_b(x)=x^2+b^2$, and let
$W(x)=\sum_{d=0}^{D}w_dx^d\in\mathbb C[x]$.  Initially in their
half-planes of absolute convergence, define
\begin{align}
 F_{W,\pm}(s)&=\sum_{n\ge0}W(x_n)(x_n\pm ib)^{-s},
 &Z_W(s)&=\sum_{n\ge0}W(x_n)Q_b(x_n)^{-s}.
 \label{sd:def:spectra}
\end{align}
The powers use principal logarithms.  Since $x_n>0$,
$\log(x_n+ib)+\log(x_n-ib)=\log Q_b(x_n)$ without a branch correction.
All evaluations at zero refer to meromorphic continuation.  We use the
normalization requested in \cite{QSpec}:
\begin{align}
 \mathscr D_W(s)&=\frac{F_{W,+}(2s)+F_{W,-}(2s)}2-Z_W(s),
 \label{sd:def:difference}\\
 \Delta_r(W;a,b)&=\mathscr D_W^{(r)}(0)
 =2^{r-1}\{F_{W,+}^{(r)}(0)+F_{W,-}^{(r)}(0)\}-Z_W^{(r)}(0).
 \label{sd:def:defect}
\end{align}
Primes and parenthesized derivatives on these spectra always
differentiate their order, not their shift or mass.

For reference, the established local anomaly is
\begin{align}
 \Delta_1(W;a,b)
 &=\frac14[x^{-1}]W(x)
       \left\{\log(1+ib/x)-\log(1-ib/x)\right\}^{2}
       \label{sd:eq:first-residue}\\
 &=\sum_{k\ge1}w_{2k-1}\frac{(-1)^kb^{2k}}{k}
       \left(H_{2k}-\frac12H_k\right).
       \label{sd:eq:first-harmonic}
\end{align}
Only finitely many terms occur.  Its independence of $a$ and its
vanishing for every $b$ precisely when $W$ is even are inherited
results, not new claims of this article.  The general local-anomaly
framework is classical; see Dowker \cite{Dowker}.  We recover the formula
below to fix the normalization and expose the additional terms at
higher orders.

\subsection{Finite harmonic coefficients and global subtraction}

For $k\ge1$ define the polynomial
\begin{equation}\label{sd:eq:Ck}
 C_k(s;b)=(-1)^kb^{2k}
 \left\{\frac{(2s)_{2k}}{(2k)!}-\frac{(s)_k}{k!}\right\}
 =\sum_{m=2}^{2k}\alpha_{k,m}(b)s^m.
\end{equation}
Here $(s)_k=s(s+1)\cdots(s+k-1)$.  For a finite list $v$,
$e_j(v)$ denotes its elementary symmetric polynomial, with $e_0=1$
and $e_j=0$ beyond the length of the list.  Set
\begin{equation}\label{sd:eq:elementary}
 E_j(N)=e_j(1,1/2,\ldots,1/N),\qquad E_j(0)=\begin{cases}1&j=0,\\0&j>0.\end{cases}
\end{equation}
These are polynomials in the generalized harmonic numbers:
\[
 \sum_{j\ge0}E_j(N)t^j
 =\exp\left(\sum_{\ell\ge1}
       \frac{(-1)^{\ell-1}}{\ell}H_N^{(\ell)}t^\ell\right).
\]

\begin{lemma}[All local coefficients]\label{sd:lem:coefficients}
For $m\ge2$,
\begin{equation}\label{sd:eq:alpha}
 \alpha_{k,m}(b)=\frac{(-1)^kb^{2k}}{k}
       \{2^{m-1}E_{m-1}(2k-1)-E_{m-1}(k-1)\}.
\end{equation}
Set $\alpha_{k,m}=0$ for $m<2$ or $m>2k$.  In particular,
\[
 C_1(s;b)=-2b^2s^2,\qquad
 \alpha_{k,2}=\frac{(-1)^kb^{2k}}k
                    (2H_{2k-1}-H_{k-1}).
\]
\end{lemma}
\begin{proof}
Use
\[
 \frac{(s)_k}{k!}=\frac{s}{k}
                    \prod_{j=1}^{k-1}(1+s/j),\qquad
 \frac{(2s)_{2k}}{(2k)!}=\frac{s}{k}
                    \prod_{j=1}^{2k-1}(1+2s/j).
\]
Their linear terms agree.  Comparing every subsequent coefficient
proves the assertion, including the stated finite support.
\end{proof}

Put
\begin{equation}\label{sd:eq:h}
 h(s,x)=\frac{(x+ib)^{-2s}+(x-ib)^{-2s}}2-Q_b(x)^{-s},
 \qquad K=\left\lfloor\frac{D+1}{2}\right\rfloor.
\end{equation}
For $D=0$ the integer $K$ is zero and all sums from $1$ to $K$ are empty.

\begin{lemma}[A globally valid analytic subtraction]
\label{sd:lem:subtraction}
In a complex neighborhood of $s=0$,
\begin{align}
 \mathscr D_W(s)
 &=\sum_{n\ge0}W(x_n)
 \left\{h(s,x_n)-x_n^{-2s}
                  \sum_{k=1}^{K}C_k(s;b)x_n^{-2k}\right\}
 \notag\\
 &\quad+\sum_{d=0}^{D}\sum_{k=1}^{K}
         w_dC_k(s;b)\zeta(2s+2k-d,a).
 \label{sd:eq:global-subtraction}
\end{align}
The infinite sum converges locally normally.  All the apparent
singularities at zero in the finite sum are removable.  The original
three spectra in \eqref{sd:def:spectra} are holomorphic at zero.
\end{lemma}
\begin{proof}
For $x>b$, the ordinary binomial expansions give
\[
 h(s,x)=x^{-2s}\sum_{k\ge1}C_k(s;b)x^{-2k}.
\]
On compact mass sets, its remainder after $K$ terms is uniformly
$O(x^{-2\Re s-2K-2})$ on a small closed $s$-disk.  After multiplication
by $W$, this is $O(x^{-2+2\varepsilon})$ when $D$ is even and
$O(x^{-3+2\varepsilon})$ when $D$ is odd, for $|s|\le\varepsilon<1/4$.
This proves normal convergence.  The finitely many $x\le b$ use the
exact expression $h$, so no assumption $b<a$ is needed.

In a common right half-plane every term being subtracted is separately
summable; adding it back as a Hurwitz zeta proves
\eqref{sd:eq:global-subtraction} there and hence by continuation.
The only possible singular terms near zero have $2k-d=1$.
The simple pole of $\zeta(1+2s,a)$ is multiplied by $C_k=O(s^2)$,
so they are removable.  Finally, expansion of the polynomial weight
about $x\pm ib$ expresses each $F_{W,\pm}$ as a finite sum of
$\zeta(s-j,a\pm ib)$ with $j\ge0$.  They are holomorphic at zero,
and the same follows for $Z_W$ from \eqref{sd:def:difference}.
\end{proof}

\subsection{The complete higher-derivative formula}

Define the elementary logarithmic polynomial
\begin{align}
 H_r(x;b)
 &=(-1)^r\left[2^{r-1}
       \{\log^r(x+ib)+\log^r(x-ib)\}-\log^r Q_b(x)\right]
       \label{sd:eq:H-log}\\
 &=\sum_{j=1}^{\lfloor r/2\rfloor}
       (-1)^{r+j}4^j\binom r{2j}
       \log^{r-2j}Q_b(x)\arctan^{2j}(b/x),
       \label{sd:eq:H-real}
\end{align}
so $H_1=0$, $H_2=-4\arctan^2(b/x)$, and
$H_3=12\log Q_b(x)\arctan^2(b/x)$.  Also put
\begin{equation}\label{sd:eq:Pr}
 P_{k,r}(L)=r!\sum_{m=2}^{\min(r,2k)}
          \alpha_{k,m}(b)\frac{(-2L)^{r-m}}{(r-m)!}.
\end{equation}

\begin{theorem}[Every order, every polynomial weight, every positive mass]
\label{sd:thm:all-orders}
For every integer $r\ge1$,
\begin{align}
 \Delta_r(W;a,b)
 &=\sum_{n\ge0}W(x_n)
       \left\{H_r(x_n;b)-\sum_{k=1}^{K}
                    \frac{P_{k,r}(\log x_n)}{x_n^{2k}}\right\}
       \notag\\
 &\quad+\sum_{d=0}^{D}\sum_{k=1}^{K}w_d\mathcal J_{k,r}(2k-d;a,b),
       \label{sd:eq:all-orders}
\end{align}
where, if $\nu\ne1$,
\begin{equation}\label{sd:eq:nonresonant}
 \mathcal J_{k,r}(\nu;a,b)=r!\sum_{m=2}^{\min(r,2k)}
   \alpha_{k,m}(b)\frac{2^{r-m}}{(r-m)!}
                       \zeta^{(r-m)}(\nu,a),
\end{equation}
and at the resonant value $\nu=1$,
\begin{equation}\label{sd:eq:resonant}
 \boxed{\displaystyle
 \mathcal J_{k,r}(1;a,b)=r!\left\{
   \frac{\alpha_{k,r+1}(b)}2+
   \sum_{m=2}^{\min(r,2k)}\alpha_{k,m}(b)
          \frac{(-2)^{r-m}}{(r-m)!}\gamma_{r-m}(a)\right\}.}
\end{equation}
The infinite sum in \eqref{sd:eq:all-orders} is ordinary and absolutely
convergent.  It and every fixed parameter derivative converge locally
normally on positive parameter sets.
\end{theorem}
\begin{proof}
Differentiate the normally convergent expression
\eqref{sd:eq:global-subtraction}.  The derivative of its exact pointwise
term is $H_r$; writing
$\log(x\pm ib)=\frac12\log Q_b(x)\pm i\arctan(b/x)$ proves
\eqref{sd:eq:H-real}.  Differentiation of $x^{-2s}C_k(s;b)$ gives
\eqref{sd:eq:Pr}.  If $\nu\ne1$, Taylor multiplication gives
\eqref{sd:eq:nonresonant}.

At the pole, use the stipulated Laurent convention:
\[
 \zeta(1+2s,a)=\frac1{2s}
       +\sum_{j\ge0}\frac{(-2)^j}{j!}\gamma_j(a)s^j.
\]
The coefficient of $s^r$ in $C_k(s;b)\zeta(1+2s,a)$ contains
$\alpha_{k,r+1}/2$ from the pole and the displayed finite convolution
from its regular part.  Multiplication by $r!$ proves
\eqref{sd:eq:resonant}.  The uniform estimates in the lemma allow
the differentiations; logarithmic factors from spectral derivatives
do not affect summability.  The same estimates on small complex
neighborhoods of positive $(a,b)$ justify parameter derivatives.
\end{proof}

At $r=1$ the ordinary sum and every nonresonant term vanish.  The
remaining term is $\alpha_{k,2}/2$, which recovers
\eqref{sd:eq:first-harmonic}.  At $r\ge2$, deleting the ordinary sum
is generally incorrect.  Deleting $\alpha_{k,r+1}/2$ is also
incorrect when this coefficient is present: differentiating a regular
part alone loses the contribution of the pole.

\subsection{A single monomial isolates every Stieltjes index}

The simplest odd weight makes the correction especially transparent.
For $r\ge2$ set
\[
 c_r=(-1)^{r-1}2^{r-1}r(r-1).
\]
Since $C_1=-2b^2s^2$, the general theorem gives
\begin{equation}\label{sd:eq:linear-weight}
 \boxed{\displaystyle
 \Delta_r(x;a,b)=
 \sum_{n\ge0}\left\{x_nH_r(x_n;b)
       -c_rb^2\frac{\log^{r-2}x_n}{x_n}\right\}
       +c_rb^2\gamma_{r-2}(a).}
\end{equation}
Its residual summands are
$O(x_n^{-3}(1+\log^{r-2}x_n))$.  Thus this is an ordinary convergent
identity for each generalized Stieltjes constant, coupled to precisely
specified spectral derivatives.  For example,
\begin{align}
 \Delta_2(x;a,b)
 &=-4\sum_{n\ge0}
       \left\{x_n\arctan^2(b/x_n)-\frac{b^2}{x_n}\right\}
       +4b^2\psi(a),\label{sd:eq:linear-second}\\
 \Delta_3(x;a,b)
 &=12\sum_{n\ge0}
       \left\{x_n\log Q_b(x_n)\arctan^2(b/x_n)
                      -\frac{2b^2\log x_n}{x_n}\right\}
       +24b^2\gamma_1(a).\label{sd:eq:linear-third}
\end{align}
At first order, $\Delta_1(x;a,b)=-b^2$ even though each ordinary
pointwise first-derivative difference is zero.  This is a useful exact
check on the chosen continuation, not a license to omit higher-order
ordinary terms.

\subsection{The first-derivative parity criterion fails completely at order two}

\begin{theorem}[Universal second-defect vanishing]
\label{sd:thm:vanishing}
Fix $a>0$.  For a polynomial $W\in\mathbb C[x]$,
\[
 \Delta_2(W;a,b)=0\ \text{for every }b>0
 \quad\Longleftrightarrow\quad W=0.
\]
Thus no nonzero even polynomial retains the universal vanishing
property of the first derivative.
\end{theorem}
\begin{proof}
For $0<b<a$, the full binomial expansion converges normally.  If
$k>(D+1)/2$, all the Hurwitz values $\zeta(2k-d,a)$ in its
$b^{2k}$ coefficient are ordinary convergent sums.  That coefficient
in $\Delta_2$ is
\begin{equation}\label{sd:eq:high-mass-coeff}
 \frac{2(-1)^k}{k}(2H_{2k-1}-H_{k-1})
             \sum_{n\ge0}\frac{W(n+a)}{(n+a)^{2k}}.
\end{equation}
The scalar is nonzero.  If the defect vanishes identically, all these
moments vanish.  If $W\ne0$, let $n_0$ be the least nonnegative integer
with $W(n_0+a)\ne0$.  Multiplying the vanishing moment by
$(n_0+a)^{2k}$ and letting $k\to\infty$ gives
$W(n_0+a)=0$, a contradiction.  Dominated convergence is valid after
fixing any sufficiently large $k_0$: polynomial growth of $W$ makes
$\sum_{n>n_0}|W(n+a)|((n_0+a)/(n+a))^{2k_0}$ finite.
The reverse implication is immediate.
\end{proof}

This is a statement about identity in the mass at a fixed shift.
It does not exclude accidental zeros at a prescribed mass for signed
polynomial weights.  No inequality or numerical independence assertion
is needed for the proof.

